% TOP_PROBLEM: {"id": 2309015, "title": "Research Problems in Function Theory — Problem 9.15", "category_id": 8, "category": {"id": 8, "name": "analysis", "display_name": "Analysis", "description": "Limits, continuity, calculus, and function theory.", "slug": "analysis", "order_index": 8, "created_at": "2026-07-31T15:26:25.671Z"}, "status": "open"}
% FIRST_POSTED: null
% ENTRY_KIND: statement_only
% Imported from: batch07.tex; UnsolvedMath ID 2309015
\documentclass[11pt]{book}
\usepackage{fontspec}
\setmainfont{Latin Modern Roman}
\setsansfont{Latin Modern Sans}
\usepackage{amsmath,amssymb,mathtools}
\usepackage{unicode-math}
\setmathfont{Latin Modern Math}
\usepackage{xurl}
\usepackage[hidelinks]{hyperref}
\newcommand{\sref}[2]{\hyperlink{source:#1:#2}{[S#2]}}
\newcommand{\cataloguescope}[1]{\paragraph{Mathematical question.}#1}
\begin{document}
% UnsolvedMath ID 2309015; source catalogue code AMR-022-9015
\setcounter{section}{2309014}
\section{Research Problems in Function Theory — Problem 9.15}\label{2309015}
{\small\sffamily UnsolvedMath ID 2309015 · Analysis}
\subsection{Definitions and mathematical statement}

\paragraph{Shared notation.}$\mathbb N=\{1,2,\ldots\}$, and $\mathbb Z,\mathbb Q,\mathbb R,\mathbb C$ denote the integers, rationals, reals and complex numbers. Any other conventions are specified locally.


Let $\mathbb D=\{z\in\mathbb C:|z|<1\}$, and let $H^{\infty}(D)$ denote the algebra of bounded analytic functions on a planar domain $D$, with supremum norm. Put $H^{\infty}=H^{\infty}(\mathbb D)$ and $I(f_1,f_2)=\{f_1h_1+f_2h_2:h_1,h_2\in H^{\infty}\}$. Let $f_1,f_2,g\in H^{\infty}$ satisfy $|g|\le|f_1|+|f_2|$ throughout $\mathbb D$.
\paragraph{Statement recorded in the supplied source.}
Is $g^2\in I(f_1,f_2)$ necessarily true? Equivalently, must $g^2=f_1h_1+f_2h_2$ for some bounded analytic $h_1,h_2$? The source records membership for the cube and counterexamples for the first power. \sref{2309015}{2}
\subsection{Short English statement}
Does modulus domination by two bounded analytic generators force the square into their ideal?
\subsection{Sources}
\begin{description}
\item[\hypertarget{source:2309015:1}{S1}] ULAM.AI and the UnsolvedMath Contributors, \emph{UnsolvedMath: A Curated Collection of Open Mathematics Problems}, record 2309015 (AMR-022-9015). CC BY 4.0. \url{https://huggingface.co/datasets/ulamai/UnsolvedMath}.
\item[\hypertarget{source:2309015:2}{S2}] Walter K. Hayman and Edward F. Lingham, \emph{Research Problems in Function Theory} (2018), arXiv:1809.07200. Problem 9.15, complete original question, all following updates, and the relevant chapter definitions. \url{https://arxiv.org/abs/1809.07200}.
\end{description}
\label{end:2309015}
\end{document}
